% ID: SEL-D0003-OA-E04
% ESTADO: cotejada. Fuente y fragmento cotejados de forma independiente; salida compilada revisada.
% Fuente: Selectividad/sel_2019_matematicas_aplicadas/Suplente Junio/Exámenes Matemáticas Aplicadas CCSS - A y B.pdf
% SHA-256: a59474e6895bd41034da53b7060969cbe8cbf7f8272f39fbf16821c95f06c84b
% Página: 1; opción A; ejercicio 4.
% Encabezado original: EJERCICIO 4. Solo enunciado; no contiene solución.

La cantidad de azúcar que añade un fabricante de refrescos a sus productos
sigue una ley Normal cuya varianza es $225\,mg^2$. Se ha seleccionado al
azar una muestra de 25 refrescos de ese fabricante, en la que se ha obtenido
una media de $175\,mg$ de azúcar añadido por refresco.

\begin{enumerate}[label=\alph*)]
  \item \textbf{(1.5 puntos)} Determine un intervalo de confianza al 90\,\%
  para la cantidad media de azúcar añadida a cada refresco.
  \item \textbf{(1 punto)} \textquestiondown{}Cuál debe ser el tamaño mínimo
  de la muestra para que el intervalo de confianza correspondiente al
  80\,\% tenga una amplitud como máximo de $5\,mg$?
\end{enumerate}
